\chapter{Environment}

Four defects stood between a clean checkout and a compiling kernel. None was
interesting in itself; all four are recorded because three of them produce
misleading rather than obvious failures.

\entry{nvcc 13.3 refuses GCC 16 as a host compiler}

\field{Symptom.} Compiling a trivial CUDA file with \texttt{-ccbin g++-16}:

\begin{lstlisting}
crt/host_config.h:137:2: error: #error -- unsupported GNU version! gcc
versions later than 15 are not supported!
\end{lstlisting}

\field{Root cause.} CUDA 13.3 supports GCC up to and including 15 and enforces
that with a preprocessor check, so this is a deliberate refusal rather than a
parse failure in some header. The machine has GCC 13, 14, 15 and 16 installed
and defaults to 15, so the problem only appears if the first design's floor of GCC
16.1 is taken literally, which it was.

\field{Options.} Pass \texttt{-allow-unsupported-compiler} and use GCC 16
anyway; drop to GCC 15.2, the newest supported host compiler; or install an
older toolkit.

\field{Fix and why.} GCC 15.2. What the floor was really buying is the C++23
host standard, and GCC 15.2 provides that in full. Suppressing the toolkit's
own compatibility check on a project whose entire output is numbers that are
supposed to be trustworthy is the wrong trade: a miscompiled force kernel that
still runs is the worst failure mode available here. Dropping the toolkit would
have given up first class \texttt{sm\_120} support on this card.

\field{Verification.} A kernel plus a CUB radix sort compiles, links and runs
correctly under \texttt{-ccbin g++-15 -arch=sm\_120 -std=c++23}.

\entry{CUDA 13 removed clockRate and memoryClockRate}

\field{Symptom.} The device probe would not compile:

\begin{lstlisting}
error: class "cudaDeviceProp" has no member "clockRate"
error: class "cudaDeviceProp" has no member "memoryClockRate"
\end{lstlisting}

\field{Root cause.} Both fields were deprecated during the CUDA 12 series and
removed in CUDA 13. The information is still available through
\texttt{cudaDeviceGetAttribute}.

\field{Fix and why.} Query both through the attribute interface. This matters
more than a compile fix normally would: the memory clock and bus width are what
the theoretical peak bandwidth is computed from, and the probe reports that figure beside
the measured bandwidth. Hard coding it from a data sheet would have meant quoting a number that had never
been measured on the machine that produced the results, which is exactly what
the no number without a run rule exists to prevent.

\field{Verification.} The probe reports \GpuBusBits{} bits at the measured
memory clock, giving \GpuPeakBandwidth{} GB/s, which agrees with the
manufacturer figure.

\entry{CMake could not identify the CUDA compiler at all}

\field{Symptom.} Configuration failed inside \texttt{project()}, before any
project code was reached:

\begin{lstlisting}
/usr/include/c++/15/type_traits(3303): error: __bool_constant is not a template
Error limit reached.
100 errors detected in the compilation of "CMakeCUDACompilerId.cu".
\end{lstlisting}

The confusing part was that nvcc had just compiled my own test kernel against
the same standard library without complaint.

\field{Root cause.} The difference was the standard flag. My test passed
\texttt{-std=c++20} explicitly. nvcc 13.3 defaults to C++17, and the libstdc++
15 headers it is paired with here use \texttt{if consteval} inside
\texttt{bits/c++config.h}, which is C++23 syntax. Under the default standard
the front end cannot parse its own standard library, so nvcc fails on an empty
translation unit. Reduced to two lines the failure is unambiguous:

\begin{lstlisting}
c++config.h(586): error: expected a "("
      if consteval { return true; } else { return false; }
\end{lstlisting}

CMake identifies the CUDA compiler by compiling a probe with
\texttt{CMAKE\_CUDA\_FLAGS} and nothing else. \texttt{CMAKE\_CUDA\_STANDARD} has
no effect there, because it applies to targets and no targets exist yet.

\field{Fix and why.} Seed \texttt{CMAKE\_CUDA\_FLAGS} with the standard before
\texttt{project()}, guarded so a value passed on the command line still wins.
Requiring every person who clones the repository to pass an extra flag would
have broken the promise that a clean clone plus two make commands reproduces
everything.

\field{Verification.} Configuration reports \texttt{The CUDA compiler
identification is NVIDIA 13.3.73 with host compiler GNU 15.2.0}.

\entry{The build silently targeted the wrong GPU architecture}

\field{Symptom.} None. Everything compiled and would have run. The only trace
was a status line reading \texttt{CUDA architectures: 75}, which is Turing, not
the Blackwell part in the machine.

\field{Root cause.} The default was being set with a guard placed after
\texttt{project()}:

\begin{lstlisting}
if(NOT DEFINED CMAKE_CUDA_ARCHITECTURES)
  set(CMAKE_CUDA_ARCHITECTURES 120)
endif()
\end{lstlisting}

Enabling the CUDA language inside \texttt{project()} assigns a default
architecture immediately, so by the time the guard runs the variable is already
defined and the branch never fires.

\field{Fix and why.} Move the guard above \texttt{project()} and make it a
cache variable, which is the only position where it can observe the undefined
state while still losing to a value passed on the command line.

The default is still \texttt{120} for development. \texttt{-DCPS\_FAT\_BINARY=ON} builds
\texttt{physics\_core} for six architectures, which is how the release and continuous
integration builds are made.

\field{Why this is the most important entry in the chapter.} Nothing broke.
The code would have run on the RTX 5070 through just in time compilation of the
embedded PTX, correctly and at the wrong speed. Every timing number in the
benchmark ladder would have been measured against a binary built for the wrong
architecture, and the only symptom would have been results somewhat worse than
expected, which is indistinguishable from the code simply being slow. A defect
that produces plausible numbers is far more dangerous than one that crashes.

% Two hyphens inside \texttt still form an en dash ligature; only a verbatim
% region suppresses it. The dash gate caught this one in the compiled PDF.
\field{Verification.} The status line now reads \texttt{CUDA architectures:
120}, and \verb|cuobjdump --list-elf| on a built object reports
\texttt{sm\_120} SASS rather than PTX alone. The latter check is the real one:
the status line reports intent, the object file reports fact.
